% \iffalse meta-comment
%
% Copyright (C) 2017- by Yanshuo Chu <yanshuoc@gmail.com>
% Copyright (C) 2025- by Wu Junhao <andrew.junhao.wu@gmail.com>
%
% This file may be distributed and/or modified under the
% conditions of the LaTeX Project Public License, either version 1.3a
% of this license or (at your option) any later version.
% The latest version of this license is in:
%
% http://www.latex-project.org/lppl.txt
%
% and version 1.3a or later is part of all distributions of LaTeX
% version 2004/10/01 or later.
%
% \fi


%
% ^^A ======================================================================
% ^^A Module internal: 自解压引导（plain TeX 跑本文件时选中 install 块）
% ^^A ======================================================================
% ^^A 同一份 dtx 走两条路：
% ^^A   tex   src/hithesis.dtx  -> \fmtname 不是 LaTeX2e，执行下面的 install 块
% ^^A   latex src/hithesis.dtx  -> \fmtname 是 LaTeX2e，跳过 install 块去排手册
% ^^A
% ^^A 跳过靠 \csname fi\endcsname：TeX 跳过条件分支时不展开 \csname，那一句
% ^^A 因此不算作 \fi，一直跳到 install 块之后那个真的 \fi 为止。
% ^^A 走 install 那条路时它展开成 \fi 收掉 \ifx，末尾的 \endbatchfile 先结束
% ^^A 运行，后面那个 \fi 读不到。
%<*internal>
\begingroup
  \def\temp{LaTeX2e}
\expandafter\endgroup\ifx\temp\fmtname\else
\csname fi\endcsname
%</internal>
%<*install>
\input docstrip

\askforoverwritefalse
\showprogress
\keepsilent

\usedir{tex/latex/hithesis}

\preamble

This is a generated file.

Copyright (C) 2017- by Chu Yanshuo <yanshuoc@gmail.com>
Copyright (C) 2025- by Wu Junhao <andrew.junhao.wu@gmail.com>

This file may be distributed and/or modified under the
conditions of the LaTeX Project Public License, either version 1.3a
of this license or (at your option) any later version.
The latest version of this license is in:

http://www.latex-project.org/lppl.txt

and version 1.3a or later is part of all distributions of LaTeX
version 2004/10/01 or later.

To produce the documentation run the original source files ending with `.dtx'
through LaTeX.

\endpreamble

\declarepreamble\cfgpreamble

This is a generated file.

Copyright (C) 2017- by Chu Yanshuo <yanshuoc@gmail.com>
Copyright (C) 2025- by Wu Junhao <andrew.junhao.wu@gmail.com>

This file may be distributed and/or modified under the
conditions of the LaTeX Project Public License, either version 1.3a
of this license or (at your option) any later version.
The latest version of this license is in:

http://www.latex-project.org/lppl.txt

and version 1.3a or later is part of all distributions of LaTeX
version 2004/10/01 or later.

This is the configuration file of the hithesis package with LaTeX2e.

\endpreamble
\declarepreamble\istpreamble

This is a generated file.

Copyright (C) 2017- by Chu Yanshuo <yanshuoc@gmail.com>
Copyright (C) 2025- by Wu Junhao <andrew.junhao.wu@gmail.com>

This file may be distributed and/or modified under the
conditions of the LaTeX Project Public License, either version 1.3a
of this license or (at your option) any later version.
The latest version of this license is in:

http://www.latex-project.org/lppl.txt

and version 1.3a or later is part of all distributions of LaTeX
version 2004/10/01 or later.

This is the configuration file of the hithesis package with LaTeX2e.

\endpreamble
\declarepreamble\docxlikepreamble

This is a generated file.

Copyright (C) 2026- by Wu Junhao <andrew.junhao.wu@gmail.com>

This file may be distributed and/or modified under the
conditions of the LaTeX Project Public License, either version 1.3c
of this license or (at your option) any later version.
The latest version of this license is in:

https://www.latex-project.org/lppl.txt

\endpreamble
\def\MetaPrefix{--}
\declarepreamble\docxlikeluapreamble

This is a generated file.

Copyright (C) 2026- by Wu Junhao <andrew.junhao.wu@gmail.com>

This file may be distributed and/or modified under the
conditions of the LaTeX Project Public License, either version 1.3c
of this license or (at your option) any later version.
The latest version of this license is in:

https://www.latex-project.org/lppl.txt

\endpreamble
\let\MetaPrefix\DoubleperCent
%</install>
%
% ^^A ---------------------------------------------------------------------
% ^^A 先把 hithesis.ins 自己生成出来。它是给拿到源码包的人用的：
% ^^A latex hithesis.ins 与 tex src/hithesis.dtx 结果一样。开发时走后者。
% ^^A ---------------------------------------------------------------------
%<*install>
\generate{\file{hithesis.ins}{\from{\jobname.dtx}{install}}}
%</install>
%
% ^^A ---------------------------------------------------------------------
% ^^A 两个真正的类。模块按运行期顺序内联进 cls，\from 的书写顺序就是它们在
% ^^A cls 里出现的顺序，与守卫块在各 dtx 里的位置无关。
% ^^A
% ^^A 三条顺序约束，改动时照着来：
% ^^A
% ^^A 一、共用文件排在每个块的最前面。docstrip 的输入文件表是全局只增的，
% ^^A     一个被多份输出共用的文件只有排在最前，各份输出要求的先后才彼此相容，
% ^^A     否则报 Incompatible order of input files。只剩一个 cls 之后这条对
% ^^A     cls 不再适用，但 hithesis.cfg 与三个旧类名的报错文件还共用 hithesis.dtx。
% ^^A     配置仍然「早读晚生效」——\hitpresetup 先记进队列，
% ^^A     cfgload 那一段调 \hit@apply@presetup 回放。
% ^^A
% ^^A 二、加载第三方宏包那一段是 *-geometry -> *-deps -> *-hyperlink 三步。
% ^^A     硬约束只有一条：hyperref 落到所有包之后，那是它的惯例，所以
% ^^A     *-hyperlink 收尾。版面参数提前是稳妥起见，让版心在任何按 \textwidth
% ^^A     算尺寸的包之前定下来；页面样式由更早的 *-pagestyle 定义，
% ^^A     不受这一步影响。deps 内部的包彼此没有约束，一块装完。
% ^^A
% ^^A 三、另外两个模块的位置有具体理由：
% ^^A     report-toc  排在 report-deps 之后，否则 enumitem 会覆盖 cft* 类长度
% ^^A     final-toc  排在最后，避免布局参数被后加载的宏包改掉
% ^^A ---------------------------------------------------------------------
%<*install>
\generate{
    \file{\jobname.cls}{
    \from{\jobname.dtx}{clshead}
    \from{hit-presetup.dtx}{shared-presetup}
    \from{hit-names.dtx}{shared-names}
    \from{hit-engine.dtx}{shared-engine}
    \from{hit-font-presets.dtx}{font-presets}
    \from{hit-messages.dtx}{shared-messages}
    \from{hit-lang.dtx}{shared-lang}
    \from{hit-fonts.dtx}{shared-fonts}
    \from{hit-calligraphy.dtx}{shared-calligraphy}
    \from{hit-math.dtx}{shared-math}
    \from{hit-table.dtx}{shared-table}
    \from{hit-footnote.dtx}{shared-footnote}
    \from{hit-list.dtx}{shared-list}
    \from{hit-hyperlink.dtx}{shared-hyperlink}
    \from{hit-title.dtx}{shared-title}
    \from{hit-spread.dtx}{shared-spread}
    \from{hit-date.dtx}{shared-date}
    \from{\jobname.dtx}{cls-docxlike}
    \from{hit-docxlike.dtx}{shared-docxlike}
    \from{hit-caption.dtx}{shared-caption}
    \from{hit-subfigure.dtx}{shared-subfigure}
    \from{hit-algorithm.dtx}{shared-algorithm}
    \from{hit-listing.dtx}{shared-listing}
    \from{hit-key-engine.dtx}{shared-key-engine}
    \from{hit-option-engine.dtx}{shared-option-engine}
    \from{hit-bibstyle.dtx}{shared-bibstyle}
    \from{hit-parts.dtx}{shared-parts}
    \from{hit-keylist.dtx}{shared-keylist}
    \from{hit-options.dtx}{shared-options}
    \from{hit-cjk.dtx}{shared-cjk}
    \from{hit-loadclass.dtx}{shared-loadclass}
    \from{hit-flow-fonts.dtx}{flow-fonts}
    \from{hit-final-pagestyle.dtx}{final-pagestyle}
    \from{hit-final-pagestyle-master.dtx}{final-pagestyle-master}
    \from{hit-final-pagestyle-bachelor.dtx}{final-pagestyle-bachelor}
    \from{hit-final-pagestyle-postdoc.dtx}{final-pagestyle-postdoc}
    \from{hit-final-pagestyle-shenzhen-bachelor.dtx}{final-pagestyle-shenzhen-bachelor}
    \from{hit-report-pagestyle.dtx}{report-pagestyle}
    \from{hit-final-cover.dtx}{final-cover}
    \from{hit-checkbox.dtx}{shared-checkbox}
    \from{hit-report-cover.dtx}{report-cover}
    \from{hit-report-cover-shenzhen-bachelor.dtx}{report-cover-shenzhen-bachelor}
    \from{hit-final-titlepage.dtx}{final-titlepage}
    \from{hit-final-abstract.dtx}{final-abstract}
    \from{hit-final-toc.dtx}{final-toc-lists}
    \from{hit-final-heading.dtx}{final-chapter}
    \from{hit-report-heading.dtx}{report-chapter}
    \from{hit-final-bib.dtx}{final-bib}
    \from{hit-report-bib.dtx}{report-bib}
    \from{hit-final-appendix.dtx}{final-appendix}
    \from{hit-final-backmatter.dtx}{final-backmatter}
    \from{hit-final-declarations.dtx}{final-declarations}
    \from{hit-final-resolution.dtx}{final-resolution}
    \from{hit-flow-geometry.dtx}{flow-geometry}
    \from{hit-flow-deps.dtx}{flow-deps}
    \from{hit-flow-hyperlink.dtx}{flow-hyperlink}
    \from{hit-final-glossary.dtx}{final-glossary}
    \from{hit-final-matter.dtx}{final-mainmatter}
    \from{hit-report-matter.dtx}{report-matter}
    \from{hit-final-paragraph.dtx}{final-paragraph}
    \from{hit-final-footnote.dtx}{final-footnote}
    \from{hit-final-math.dtx}{final-math}
    \from{hit-final-float.dtx}{final-floats}
    \from{hit-report-float.dtx}{report-floats}
    \from{hit-report-subfig.dtx}{report-subfig}
    \from{hit-final-toc.dtx}{final-toc}
    \from{hit-report-toc.dtx}{report-toc}
    \from{hit-numbering.dtx}{shared-numbering}
    \from{hit-bilingual.dtx}{shared-bilingual}
    \from{\jobname.dtx}{clstail}
    \from{\jobname.dtx}{cfgload}
    }
    \usepreamble\cfgpreamble
    \file{\jobname.cfg}{
    \from{\jobname.dtx}{hithesis-cfg}
    \from{hit-info.dtx}{hithesis-cfg}
    \from{hit-terms.dtx}{hithesis-cfg}
    \from{hit-layout.dtx}{hithesis-cfg}
    \from{hit-thm.dtx}{hithesis-cfg}
    }
    \usepreamble\istpreamble
    \file{\jobname.ist}{\from{hit-ist.dtx}{ist}}
    \usepreamble\defaultpreamble\usepostamble\defaultpostamble
    \file{\jobname book.cls}{\from{\jobname.dtx}{book-stub}}
    \file{\jobname art.cls}{\from{\jobname.dtx}{art-stub}}
    \file{\jobname artplus.cls}{\from{\jobname.dtx}{artplus-stub}}
    \file{\jobname.sty}{\from{hit-sty.dtx}{hit-sty}}
    \usepreamble\docxlikepreamble
    \file{docxlike.sty}{
    \from{docxlike.dtx}{package}
    \from{docxlike-page.dtx}{docxlike-page}
    \from{docxlike-metrics.dtx}{docxlike-metrics}
    \from{docxlike-format.dtx}{docxlike-format}
    \from{docxlike-linebreak.dtx}{docxlike-linebreak}
    \from{docxlike-headfoot.dtx}{docxlike-headfoot}
    \from{docxlike-shape.dtx}{docxlike-shape}
    \from{docxlike-note.dtx}{docxlike-note}
    \from{docxlike-table.dtx}{docxlike-table}
    \from{docxlike-numbering.dtx}{docxlike-numbering}
    \from{docxlike-heading.dtx}{docxlike-heading}
    \from{docxlike-preset.dtx}{docxlike-preset}
    }
    \usepreamble\defaultpreamble
    \file{dtx-style.sty}{\from{hit-dtx-style.dtx}{dtx-style}}
    \file{\jobname-calligraphy.def}{\from{hit-calligraphy-data.dtx}{shared-calligraphy-data}}
}
%</install>
%
% ^^A ---------------------------------------------------------------------
% ^^A bst 不要 postamble，docstrip 的 %% 结尾行 bibtex 不认。
% ^^A docxlike 的四个生成物用它自己的版权声明（\docxlikepreamble，与 latex-docxlike 的
% ^^A docxlike.ins 同一份）；两个 .lua 的注释前缀是 --，声明前言与生成时都得换过去。
% ^^A ---------------------------------------------------------------------
%<*install>
\generate{
    \usepostamble\empty
    \file{\jobname.bst}{\from{hit-bst.dtx}{bst}}
}

\def\MetaPrefix{--}
\generate{
    \usepreamble\docxlikeluapreamble\nopostamble
    \file{docxlike.lua}{\from{docxlike-linebreak.dtx}{lua}}
    \file{docxlike-run.lua}{\from{docxlike-format.dtx}{run-lua}}
    \file{docxlike-font-metrics.lua}{\from{docxlike-metrics.dtx}{metrics-lua}}
}
\let\MetaPrefix\DoubleperCent
\generate{
    \usepreamble\docxlikepreamble\nopostamble
    \file{docxlike-font-metrics.def}{\from{docxlike-metrics.dtx}{metrics-table}}
}

\generate{\usepreamble\defaultpreamble\nopostamble
    \file{aerostatic.tex}{\from{hit-example-figures.dtx}{figure-common,aerostatic,figure-end}}
    \file{aerodynamic.tex}{\from{hit-example-figures.dtx}{figure-common,aerodynamic,figure-end}}
    \file{hybrid.tex}{\from{hit-example-figures.dtx}{figure-common,hybrid,figure-end}}
    \file{squeeze-film.tex}{\from{hit-example-figures.dtx}{figure-common,squeeze-film,figure-end}}
    \file{volute-section.tex}{\from{hit-example-figures.dtx}{volute-section}}
}

\ifToplevel{%
  \Msg{***********************************************************}
  \Msg{*}
  \Msg{* To finish the installation you have to move the following}
  \Msg{* files into a directory searched by TeX:}
  \Msg{*}
  \Msg{* The recommended directory is TEXMF/tex/latex/hithesis}
  \Msg{*}
  \Msg{* \space\space hithesis.cls}
  \Msg{* \space\space hithesisbook.cls}
  \Msg{* \space\space hithesisart.cls}
  \Msg{* \space\space hithesisartplus.cls}
  \Msg{* \space\space hithesis.sty}
  \Msg{* \space\space docxlike.sty}
  \Msg{* \space\space hithesis.ist}
  \Msg{* \space\space hithesis.cfg}
  \Msg{* \space\space hithesis.bst}
  \Msg{* \space\space docxlike.lua}
  \Msg{* \space\space docxlike-run.lua}
  \Msg{* \space\space docxlike-font-metrics.lua}
  \Msg{* \space\space docxlike-font-metrics.def}
  \Msg{* \space\space hithesis-calligraphy.def}
  \Msg{*}
  \Msg{* To produce the documentation run the files ending with}
  \Msg{* `.dtx' through LaTeX.}
  \Msg{*}
  \Msg{* Happy TeXing!}
  \Msg{***********************************************************}}

\endbatchfile
%</install>
%<*internal>
\fi
%</internal>
%
% \iffalse
% ^^A ======================================================================
% ^^A Module driver: 手册渲染驱动（latex hithesis.dtx 时由它驱动）
% ^^A ======================================================================
%<*driver>
\ProvidesFile{hithesis.dtx}[0000/00/00 4.0a Harbin Institute of Technology Thesis Template]
% 手册由类自己排，博士学位论文那一档。字体：装了中易宋体就按 windows 那套排
% （学校要的字体），没装（CI 的 runner）退到随 TeX Live 发的 fandol。
% \IfFontExistsTF 要 fontspec，这时候还没有，用 \font 探
\suppressfontnotfounderror=1
\font\hitmanualprobe="SimSun"\relax
\ifx\hitmanualprobe\nullfont
  \PassOptionsToClass{fontset=fandol}{hithesis}
\else
  \PassOptionsToClass{fontset=windows}{hithesis}
\fi
\suppressfontnotfounderror=0
% 索引的封装由 hypdoc 加，hyperref 的 hyperindex 得关着（见 dtx-style）
\PassOptionsToPackage{hyperindex=false}{hyperref}
\documentclass[stage=final,degree-level=doctor,campus=harbin]{hithesis}
\usepackage{dtx-style}

\ifdefined\hitDeveloperManual
  \EnableCrossrefs
  \CodelineIndex
\fi
\RecordChanges

\begin{document}
  % 封面、摘要、目录、正文三章都在手册文件里
  \DocInput{src/manual/hit-manual.dtx}
  \ifdefined\hitDeveloperManual
  % 实现全部进附录；修改记录与索引接在附录后，最后是致谢
  \backmatter
  \begin{appendix}
  % 给读者看的几章（写作与排版建议、软件推荐、课题组招聘）在实现之前
  \DocInput{src/manual/hit-manual-appendix.dtx}
  \chapter{实现}[Implementation]
  % 写死文件名不写 \jobname：换个 -jobname 编时 \jobname 指的是输出名，
  % 不是源文件名
  \DocInput{src/hithesis.dtx}
  \DocInput{src/hit-bst.dtx}
  \DocInput{src/hit-ist.dtx}
  \DocInput{src/hit-dtx-style.dtx}
  \DocInput{src/hit-sty.dtx}
  \DocInput{src/engine/hit-options.dtx}
  \DocInput{src/engine/hit-cjk.dtx}
  \DocInput{src/engine/hit-loadclass.dtx}
  \DocInput{src/flow/hit-flow-deps.dtx}
  \DocInput{src/flow/hit-final-glossary.dtx}
  \DocInput{src/flow/hit-flow-fonts.dtx}
  \DocInput{src/flow/hit-flow-hyperlink.dtx}
  \DocInput{src/flow/hit-final-footnote.dtx}
  \DocInput{src/flow/hit-final-math.dtx}
  \DocInput{src/flow/hit-report-subfig.dtx}
  \DocInput{src/engine/hit-keylist.dtx}
  \DocInput{src/pages/hit-final-cover.dtx}
  \DocInput{src/pages/hit-final-titlepage.dtx}
  \DocInput{src/pages/hit-final-abstract.dtx}
  \DocInput{src/pages/hit-final-toc.dtx}
  \DocInput{src/pages/hit-final-appendix.dtx}
  \DocInput{src/pages/hit-final-backmatter.dtx}
  \DocInput{src/pages/hit-final-declarations.dtx}
  \DocInput{src/pages/hit-final-resolution.dtx}
  \DocInput{src/pages/hit-final-bib.dtx}
  \DocInput{src/pages/hit-checkbox.dtx}
  \DocInput{src/pages/hit-report-cover.dtx}
  \DocInput{src/pages/hit-report-cover-shenzhen-bachelor.dtx}
  \DocInput{src/pages/hit-report-toc.dtx}
  \DocInput{src/pages/hit-report-bib.dtx}
  \DocInput{src/flow/hit-flow-geometry.dtx}
  \DocInput{src/flow/hit-final-pagestyle.dtx}
  \DocInput{src/flow/hit-final-pagestyle-master.dtx}
  \DocInput{src/flow/hit-final-pagestyle-bachelor.dtx}
  \DocInput{src/flow/hit-final-pagestyle-postdoc.dtx}
  \DocInput{src/flow/hit-final-pagestyle-shenzhen-bachelor.dtx}
  \DocInput{src/flow/hit-final-heading.dtx}
  \DocInput{src/flow/hit-final-paragraph.dtx}
  \DocInput{src/flow/hit-final-float.dtx}
  \DocInput{src/flow/hit-final-matter.dtx}
  \DocInput{src/flow/hit-report-pagestyle.dtx}
  \DocInput{src/flow/hit-report-heading.dtx}
  \DocInput{src/flow/hit-report-float.dtx}
  \DocInput{src/flow/hit-report-matter.dtx}
  \DocInput{src/engine/hit-presetup.dtx}
  \DocInput{src/config/hit-names.dtx}
  \DocInput{src/config/hit-font-presets.dtx}
  \DocInput{src/config/hit-info.dtx}
  \DocInput{src/config/hit-terms.dtx}
  \DocInput{src/config/hit-layout.dtx}
  \DocInput{src/config/hit-thm.dtx}
  \DocInput{src/engine/hit-engine.dtx}
  \DocInput{src/engine/hit-messages.dtx}
  \DocInput{src/engine/hit-lang.dtx}
  \DocInput{src/engine/hit-option-engine.dtx}
  \DocInput{src/engine/hit-key-engine.dtx}
  \DocInput{src/engine/hit-fonts.dtx}
  \DocInput{src/engine/hit-calligraphy.dtx}
  \DocInput{src/engine/hit-math.dtx}
  \DocInput{src/engine/hit-table.dtx}
  \DocInput{src/engine/hit-list.dtx}
  \DocInput{src/engine/hit-footnote.dtx}
  \DocInput{src/engine/hit-hyperlink.dtx}
  \DocInput{src/engine/hit-title.dtx}
  \DocInput{src/engine/hit-spread.dtx}
  \DocInput{src/engine/hit-date.dtx}
  \DocInput{src/docxlike/docxlike.dtx}
  \DocInput{src/docxlike/docxlike-page.dtx}
  \DocInput{src/docxlike/docxlike-metrics.dtx}
  \DocInput{src/docxlike/docxlike-format.dtx}
  \DocInput{src/docxlike/docxlike-linebreak.dtx}
  \DocInput{src/docxlike/docxlike-headfoot.dtx}
  \DocInput{src/docxlike/docxlike-shape.dtx}
  \DocInput{src/docxlike/docxlike-note.dtx}
  \DocInput{src/docxlike/docxlike-table.dtx}
  \DocInput{src/docxlike/docxlike-numbering.dtx}
  \DocInput{src/docxlike/docxlike-heading.dtx}
  \DocInput{src/docxlike/docxlike-preset.dtx}
  \DocInput{src/engine/hit-docxlike.dtx}
  \DocInput{src/engine/hit-caption.dtx}
  \DocInput{src/engine/hit-subfigure.dtx}
  \DocInput{src/engine/hit-algorithm.dtx}
  \DocInput{src/engine/hit-listing.dtx}
  \DocInput{src/engine/hit-numbering.dtx}
  \DocInput{src/engine/hit-bilingual.dtx}
  \DocInput{src/engine/hit-bibstyle.dtx}
  \DocInput{src/engine/hit-parts.dtx}
  \end{appendix}
  \else
  \input{hithesis-changes.tex}
  \fi
  \makeatletter
  \hit@doc@chapter{更新日志}{Change History}
  \PrintChanges
  \ifdefined\hitDeveloperManual
  \hit@doc@chapter{索引}{Index}
  \PrintIndex
  \fi
  \makeatother
  \DocInput{src/manual/hit-manual-thanks.dtx}
\end{document}
%</driver>
% \fi
%
% \CheckSum{0}
%
% \CharacterTable
%  {Upper-case    \A\B\C\D\E\F\G\H\I\J\K\L\M\N\O\P\Q\R\S\T\U\V\W\X\Y\Z
%   Lower-case    \a\b\c\d\e\f\g\h\i\j\k\l\m\n\o\p\q\r\s\t\u\v\w\x\y\z
%   Digits        \0\1\2\3\4\5\6\7\8\9
%   Exclamation   \!     Double quote  \"     Hash (number) \#
%   Dollar        \$     Percent       \%     Ampersand     \&
%   Acute accent  \'     Left paren    \(     Right paren   \)
%   Asterisk      \*     Plus          \+     Comma         \,
%   Minus         \-     Point         \.     Solidus       \/
%   Colon         \:     Semicolon     \;     Less than     \<
%   Equals        \=     Greater than  \>     Question mark \?
%   Commercial at \@     Left bracket  \[     Backslash     \\
%   Right bracket \]     Circumflex    \^     Underscore    \_
%   Grave accent  \`     Left brace    \{     Vertical bar  \|
%   Right brace   \}     Tilde         \~}
%
% \DoNotIndex{\newenvironment,\@bsphack,\@empty,\@esphack,\sfcode}
% \DoNotIndex{\addtocounter,\label,\let,\linewidth,\newcounter}
% \DoNotIndex{\noindent,\normalfont,\par,\parskip,\phantomsection}
% \DoNotIndex{\providecommand,\ProvidesPackage,\refstepcounter}
% \DoNotIndex{\RequirePackage,\setcounter,\setlength,\string,\strut}
% \DoNotIndex{\textbackslash,\texttt,\ttfamily,\usepackage}
% \DoNotIndex{\begin,\end,\begingroup,\endgroup,\par,\\}
% \DoNotIndex{\if,\ifx,\ifdim,\ifnum,\ifcase,\else,\or,\fi}
% \DoNotIndex{\let,\def,\xdef,\edef,\newcommand,\renewcommand}
% \DoNotIndex{\expandafter,\csname,\endcsname,\relax,\protect}
% \DoNotIndex{\Huge,\huge,\LARGE,\Large,\large,\normalsize}
% \DoNotIndex{\small,\footnotesize,\scriptsize,\tiny}
% \DoNotIndex{\centering,\raggedright,\ref}
% \DoNotIndex{\c@secnumdepth,\@startsection,\@setfontsize}
% \DoNotIndex{\ ,\@plus,\@minus,\p@,\z@,\@m,\@M,\@ne,\m@ne}
% \DoNotIndex{\@@par,\DeclareOperation,\RequirePackage,\LoadClass}
% \DoNotIndex{\AtBeginDocument,\AtEndDocument}
%
% ^^A 修改记录与索引由 driver 在附录之后排，这里不再借 \Finale 印
% \StopEventually{}
% \clearpage
%
% \section{装配表}
%
% \changes{v4.0a}{2026/09/13}{内核下限为 2023-06-01，与 wiseref 的要求一致；
%   类改用 \cs{ProvidesExplClass} 常开 expl3 语法；只剩一个 \file{hithesis.cls}，
%   终稿与报告的差别由 \option{stage} 分岔，两条装配线的配置回放合成一段
%   \texttt{cfgload}，三个旧类名只留一份报错的类文件；索引排序样式与手册样式各自成文件}
%
% ^^A ======================================================================
% ^^A 一个类由下面这些文件按这个顺序拼出来。权威清单在本文件的 install 守卫里，
% ^^A 这张表是它的产物，改完那边跑 make mapfix。
% ^^A \from 的书写顺序就是各块在 cls 里出现的顺序，与守卫在各 dtx 里的位置无关。
% ^^A ======================================================================
% ^^A ---8<--- 装配表由 scripts/check-map.py --fix 生成，别手改 ---8<---
% ^^A ======================================================================
% ^^A hithesis 的装配表
% ^^A
% ^^A 这个类由下面这些文件按这个顺序拼出来。权威清单在 src/hithesis.dtx
% ^^A 的 install 守卫里，这张表是它的产物，改完那边跑 make mapfix。
% ^^A \from 的书写顺序就是各块在 cls 里出现的顺序，与守卫在各 dtx 里的
% ^^A 位置无关。★ 标的几块在本文件里，其余在 setup/ utils/ flow/ pages/ 下。
% ^^A ======================================================================
% ^^A  1  hithesis.dtx                      clshead                  类文件头：\NeedsTeXFormat、\ProvidesExplClass
% ^^A  2  engine/hit-presetup.dtx           shared-presetup          配置队列的机制：记一条、挂条件、回放
% ^^A  3  config/hit-names.dtx              shared-names             编号后间距与模板固定名字
% ^^A  4  engine/hit-engine.dtx             shared-engine            引擎差异的封装：xetex 与 luatex 各一套
% ^^A  5  config/hit-font-presets.dtx       font-presets             字体预设、角色继承与显式覆盖
% ^^A  6  engine/hit-messages.dtx           shared-messages          面向用户的提示文案
% ^^A  7  engine/hit-lang.dtx               shared-lang              支持哪几种语言
% ^^A  8  engine/hit-fonts.dtx              shared-fonts             字号命令与字体族
% ^^A  9  engine/hit-calligraphy.dtx        shared-calligraphy       照轮廓画的校名字标与美术字
% ^^A 10  engine/hit-math.dtx               shared-math              数学与定理环境
% ^^A 11  engine/hit-table.dtx              shared-table             长表格字号
% ^^A 12  engine/hit-footnote.dtx           shared-footnote          脚注格式
% ^^A 13  engine/hit-list.dtx               shared-list              列表间距与 publist 环境
% ^^A 14  engine/hit-hyperlink.dtx          shared-hyperlink         超链接基础设施
% ^^A 15  engine/hit-title.dtx              shared-title             题目的分行与副标题
% ^^A 16  engine/hit-spread.dtx             shared-spread            定宽拉伸排字
% ^^A 17  engine/hit-date.dtx               shared-date              中文日期格式
% ^^A 18  hithesis.dtx                      cls-docxlike             加载排版引擎 docxlike：Word 页面设置、段落格式、断行
% ^^A 19  engine/hit-docxlike.dtx           shared-docxlike          类接到 docxlike 上：题序后间距、兼容桥、开关同步
% ^^A 20  engine/hit-caption.dtx            shared-caption           题注与双语题注
% ^^A 21  engine/hit-subfigure.dtx          shared-subfigure         子图
% ^^A 22  engine/hit-algorithm.dtx          shared-algorithm         伪代码：algorithm2e 的外框、关键字、行号预设
% ^^A 23  engine/hit-listing.dtx            shared-listing           代码清单：listings 的外框、行号、题注预设
% ^^A 24  engine/hit-key-engine.dtx         shared-key-engine        键怎么造、怎么存、旧名怎么兼容
% ^^A 25  engine/hit-option-engine.dtx      shared-option-engine     类选项怎么声明、标志位怎么读
% ^^A 26  engine/hit-bibstyle.dtx           shared-bibstyle          参考文献样式的开关
% ^^A 27  engine/hit-parts.dtx              shared-parts             部件排不排、只排一次、前置正文后置的切换、\hitif
% ^^A 28  engine/hit-keylist.dtx            shared-keylist           有哪些信息字段与常量
% ^^A 29  engine/hit-options.dtx            shared-options           文档类选项：声明、默认值、\ProcessKeyOptions 之后的后处理
% ^^A 30  engine/hit-cjk.dtx                shared-cjk               不经 ctex 的中文层：字号、\ccwd、LuaTeX 的中文字族
% ^^A 31  engine/hit-loadclass.dtx          shared-loadclass         加载 book 或 ctexbook 基础类
% ^^A 32  flow/hit-flow-fonts.dtx           flow-fonts               正文字号与主字体
% ^^A 33  flow/hit-final-pagestyle.dtx      final-pagestyle          页眉页脚
% ^^A 34  flow/hit-final-pagestyle-master.dtx final-pagestyle-master   页眉的硕士变体
% ^^A 35  flow/hit-final-pagestyle-bachelor.dtx final-pagestyle-bachelor 页眉的本科变体
% ^^A 36  flow/hit-final-pagestyle-postdoc.dtx final-pagestyle-postdoc  页眉的博士后变体
% ^^A 37  flow/hit-final-pagestyle-shenzhen-bachelor.dtx final-pagestyle-shenzhen-bachelor 页眉的深圳本科变体
% ^^A 38  flow/hit-report-pagestyle.dtx     report-pagestyle         页眉页脚
% ^^A 39  pages/hit-final-cover.dtx         final-cover              封面（封一）
% ^^A 40  pages/hit-checkbox.dtx            shared-checkbox          勾选框
% ^^A 41  pages/hit-report-cover.dtx        report-cover             封面（内含目录）
% ^^A 42  pages/hit-report-cover-shenzhen-bachelor.dtx report-cover-shenzhen-bachelor 深圳本科报告首页的变体
% ^^A 43  pages/hit-final-titlepage.dtx     final-titlepage          内封（扉页）
% ^^A 44  pages/hit-final-abstract.dtx      final-abstract           摘要
% ^^A 45  pages/hit-final-toc.dtx           final-toc-lists          图表目录
% ^^A 46  flow/hit-final-heading.dtx        final-chapter            章节标题
% ^^A 47  flow/hit-report-heading.dtx       report-chapter           章节标题
% ^^A 48  pages/hit-final-bib.dtx           final-bib                参考文献版式
% ^^A 49  pages/hit-report-bib.dtx          report-bib               参考文献版式
% ^^A 50  pages/hit-final-appendix.dtx      final-appendix           附录
% ^^A 51  pages/hit-final-backmatter.dtx    final-backmatter         成果、索引、简历
% ^^A 52  pages/hit-final-declarations.dtx  final-declarations       原创性声明与使用授权
% ^^A 53  pages/hit-final-resolution.dtx    final-resolution         答辩决议
% ^^A 54  flow/hit-flow-geometry.dtx        flow-geometry            版心，提前定下来
% ^^A 55  flow/hit-flow-deps.dtx            flow-deps                多模块共用的第三方宏包
% ^^A 56  flow/hit-flow-hyperlink.dtx       flow-hyperlink           hyperref 配置，排在所有包之后
% ^^A 57  flow/hit-final-glossary.dtx       final-glossary           术语与符号表
% ^^A 58  flow/hit-final-matter.dtx         final-mainmatter         骨架宏与前后文切换
% ^^A 59  flow/hit-report-matter.dtx        report-matter            骨架宏
% ^^A 60  flow/hit-final-paragraph.dtx      final-paragraph          段落
% ^^A 61  flow/hit-final-footnote.dtx       final-footnote           脚注
% ^^A 62  flow/hit-final-math.dtx           final-math               公式断行
% ^^A 63  flow/hit-final-float.dtx          final-floats             浮动体
% ^^A 64  flow/hit-report-float.dtx         report-floats            浮动体
% ^^A 65  flow/hit-report-subfig.dtx        report-subfig            子图兼容层
% ^^A 66  pages/hit-final-toc.dtx           final-toc                目录版式
% ^^A 67  pages/hit-report-toc.dtx          report-toc               目录版式
% ^^A 68  engine/hit-numbering.dtx          shared-numbering         图表与公式编号的重置关系和格式
% ^^A 69  engine/hit-bilingual.dtx          shared-bilingual         主副语言内容解析与输入方向
% ^^A 70  hithesis.dtx                      clstail                  \AtEndOfClass 钩子
% ^^A 71  hithesis.dtx                      cfgload                  读 hithesis.cfg 并回放配置队列
% ^^A ======================================================================
% ^^A ---8<--- 装配表到此 ---8<---
%
% \section{实现细节}
%
% \subsection{基本信息}
%
%    \begin{macrocode}
%<clshead>\NeedsTeXFormat{LaTeX2e}[2023/06/01]
%<clshead>\ProvidesExplClass{hithesis}{0000/00/00}{4.0a}{Harbin Institute of Technology Thesis Template}
%<*hithesis-cfg>
\ExplSyntaxOff
\ProvidesExplFile{hithesis.cfg}{0000/00/00}{4.0a}{Harbin Institute of Technology Thesis Template}
%</hithesis-cfg>
%    \end{macrocode}
%
% \subsection{排版引擎}
%
% 以 preset=none 加载 docxlike，由文档类提供默认版式。
% src/docxlike 是独立仓库的副本，修改正本后用 make docxlike-sync 同步。
%    \begin{macrocode}
%<cls-docxlike>\RequirePackage[preset=none]{docxlike}
%<cls-docxlike>\NewCommandCopy \hitenter \parmark
%    \end{macrocode}
%
% \subsection{按阶段初始化模块}
% \label{sec:stage-gate}
%
% 阶段条件只选择执行分支，不引入 TeX 分组。
% 条件参数已完成记号化，不能在其内部用 ExplSyntaxOff 改变自身的 catcode。
%
% \subsection{其它}
% \label{sec:other}
% 在模板文档结束时即装入配置文件，这样用户就能在导言区进行相应的修改。
%
% \iffalse
% ^^A 示例插图源码位于 src/assets/hit-example-figures.dtx，docstrip 提取后由 make figures 编译。
% \fi
%
% \section{类尾与配置回放}
%
% 所有模块加载后回放默认配置，再解析词条和部件变体。
%
% 类末尾开 \cs{sloppy}：Word 里的行从不伸出版心，排不下就把词隙拉开。交给 \hologo{TeX}
% 断行的段（\hologo{XeTeX} 与 \hologo{pdfTeX} 下的正文，\hologo{LuaTeX} 下 docxlike 不接手的段）
% 照样宁松勿溢。它不动 \cs{hbadness}，松的行照报 Underfull，只有溢出 0.5\,pt 以内的不报。
%
%    \begin{macrocode}
%<*clstail>
\AtEndOfClass{\sloppy}
\hit@setup@numbering
%</clstail>
%    \end{macrocode}
%
% 词条解析必须紧跟默认配置回放，不能推迟到导言区之后，否则会覆盖用户设置。
%
% 页眉页脚的字体只有正文照终稿排的那几档有目标。
%    \begin{macrocode}
%<*cfgload>
\input{hithesis.cfg}
\hit@axis@bools:
\hit@apply@presetup
\hit@axis@resolve:
\hit@hook@resolve:
\AfterPreamble
  {
    \hit@docxlike@apply@body:
    \bool_if:NT \g__hit_final_body_bool
      {
        \docxlike_format_apply_hf:nN { Header } \hit@header@font
        \docxlike_format_apply_hf:nN { Footer } \hit@footer@font
      }
  }
%</cfgload>
%    \end{macrocode}
%
% \subsection{三个旧类名}
% 三个旧类名计划在 v4 移除，预计约 2027 年初发布。稳定版 v3.2a 继续保留原有类名。
% 当前为 v4 开发版，只保留迁移错误提示，指向统一的 hithesis 类后立即停止。
%
% \begin{longtblr}[caption={旧类名改成什么}, label={tab:impl-oldclass}]
%   {colspec={l X[2,l]}, width=\linewidth}
% \toprule
% 旧写法 & 改成 \\ \midrule
% \cs{documentclass}\marg{hithesisbook} &
%   \cs{documentclass}\oarg{stage=final,\ldots}\marg{hithesis} \\
% \cs{documentclass}\marg{hithesisart} &
%   \cs{documentclass}\oarg{stage=proposal,\ldots}\marg{hithesis} \\
% \cs{documentclass}\marg{hithesisartplus} &
%   \cs{documentclass}\oarg{stage=interim,degree-level=doctor,campus=shenzhen}\marg{hithesis} \\
% \bottomrule
% \end{longtblr}
%
% \subsection{\file{hithesisbook.cls}}
%    \begin{macrocode}
%<*book-stub>
\NeedsTeXFormat{LaTeX2e}[2023/06/01]
\ProvidesClass{hithesisbook}
  [0000/00/00 v4.0a hithesisbook removed, see error]
\ClassError{hithesisbook}
  {%
    现已不再支持调用 hithesisbook 类，请改为 hithesis%
  }
  {%
    把 \string\documentclass 那一行改成：\MessageBreak
    \space\space\string\documentclass[stage=final, ...]{hithesis}\MessageBreak
    \MessageBreak
    三个旧类名在 v4 移除，预计约 2027 年初发布。\MessageBreak
    稳定版 v3.2a 仍保留原有类名；当前为 v4 开发版。\MessageBreak
    它们原先只是把选项转给真正的类，而现在统一从 hithesis 进，\MessageBreak
    文档类型由 stage 选择：final 是学位论文，proposal 与 interim 是\MessageBreak
    开题与中期报告。\MessageBreak
    文档其余部分一个字都不用改%
  }
\stop
%</book-stub>
%    \end{macrocode}
%
% \subsection{\file{hithesisart.cls}}
%    \begin{macrocode}
%<*art-stub>
\NeedsTeXFormat{LaTeX2e}[2023/06/01]
\ProvidesClass{hithesisart}
  [0000/00/00 v4.0a hithesisart removed, see error]
\ClassError{hithesisart}
  {%
    类 hithesisart 已移除，请改用 hithesis%
  }
  {%
    把 \string\documentclass 那一行改成：\MessageBreak
    \space\space\string\documentclass[stage=proposal, ...]{hithesis}\space\space(或 stage=interim)\MessageBreak
    \MessageBreak
    三个旧类名在 v4 移除，预计约 2027 年初发布。\MessageBreak
    稳定版 v3.2a 仍保留原有类名；当前为 v4 开发版。\MessageBreak
    它们原先只是把选项转给真正的类，而现在统一从 hithesis 进，\MessageBreak
    文档类型由 stage 选择：final 是学位论文，proposal 与 interim 是\MessageBreak
    开题与中期报告。\MessageBreak
    文档其余部分一个字都不用改%
  }
\stop
%</art-stub>
%    \end{macrocode}
%
% \subsection{\file{hithesisartplus.cls}}
%    \begin{macrocode}
%<*artplus-stub>
\NeedsTeXFormat{LaTeX2e}[2023/06/01]
\ProvidesClass{hithesisartplus}
  [0000/00/00 v4.0a hithesisartplus removed, see error]
\ClassError{hithesisartplus}
  {%
    类 hithesisartplus 已移除，请改用 hithesis%
  }
  {%
    把 \string\documentclass 那一行改成：\MessageBreak
    \space\space\string\documentclass[stage=interim,degree-level=doctor,campus=shenzhen]{hithesis}\MessageBreak
    \MessageBreak
    三个旧类名在 v4 移除，预计约 2027 年初发布。\MessageBreak
    稳定版 v3.2a 仍保留原有类名；当前为 v4 开发版。\MessageBreak
    它们原先只是把选项转给真正的类，而现在统一从 hithesis 进，\MessageBreak
    文档类型由 stage 选择：final 是学位论文，proposal 与 interim 是\MessageBreak
    开题与中期报告。\MessageBreak
    文档其余部分一个字都不用改%
  }
\stop
%</artplus-stub>
%    \end{macrocode}
%
% \Finale
%
\endinput
% \iffalse
%  Local Variables:
%  mode: doctex
%  TeX-master: t
%  End:
% \fi
